\documentclass[11pt,a4paper]{article}

% ---------------------------------------------------------
% PREAMBLE (stable, warning-free, Overleaf-safe)
% ---------------------------------------------------------
\usepackage[utf8]{inputenc}
\usepackage[T1]{fontenc}
\usepackage{lmodern}
\usepackage{amsmath, amssymb, amsfonts}
\usepackage{bm}
\usepackage{geometry}
\usepackage{graphicx}
\usepackage{hyperref}
\usepackage{cite}
\usepackage{authblk}
\usepackage{physics}
\usepackage{mathtools}

\geometry{margin=2.4cm}

\title{\textbf{Companion XXXVII: Disk Dynamics and Rotation Curves in the Relaxation-Driven Cyclic Cosmology}}
\author{Michael Lehmann \\ \texttt{mi.lehmann@gmx.de}}
\date{April 2026}

\begin{document}
\maketitle

\begin{abstract}
The rotation curves of disk galaxies provide one of the most direct observational 
probes of gravitational dynamics on kiloparsec scales. Their characteristic flatness 
is traditionally interpreted as evidence for extended dark matter halos. In the 
Relaxation-Driven Cyclic Cosmology (RDCC), the scale-dependent evolution of the 
gravitational potential modifies the internal structure of halos and the distribution 
of angular momentum in a way that reshapes the dynamics of galactic disks. This 
Companion develops a continuous description of disk formation, derives the RDCC 
rotation curve, and explains how the relaxation scale influences the transition 
between rising, flat, and declining regimes. The resulting framework provides a 
natural interpretation of observed rotation curves without invoking exotic dark 
matter microphysics.
\end{abstract}
% ---------------------------------------------------------
\section{Introduction}

Disk galaxies are among the most dynamically ordered structures in the Universe. 
Their rotation curves, which remain approximately flat far beyond the luminous 
extent of the disk, have long been regarded as a cornerstone of the dark matter 
paradigm. In the standard cosmological model, this flatness arises from the 
presence of an extended halo whose density profile declines more slowly than the 
baryonic component.

In RDCC, the gravitational potential evolves differently. The relaxation mechanism 
suppresses the growth of small-scale modes and enhances the coherence of large-scale 
modes. This modifies the density structure of halos, the distribution of angular 
momentum, and the gravitational field experienced by the disk. The rotation curve 
becomes a direct tracer of the relaxation scale $k_{\mathrm{rel}}$ and the temporal 
structure of the potential.
% ---------------------------------------------------------
\section{The RDCC Gravitational Potential and Disk Formation}

The formation of a galactic disk begins with the collapse of gas within a dark 
matter halo. The gravitational potential governing this collapse is modified in 
RDCC according to
\begin{equation}
    \Phi_{\mathrm{RDCC}}(k,a)
    = \Phi(k,a)
      \left[1 - \chi_{\Phi}
      \left(\frac{k}{k_{\mathrm{rel}}}\right)^{\alpha}\right].
\end{equation}

This modification affects the radial distribution of mass within the halo. 
Small-scale modes, which would normally steepen the inner density profile, are 
suppressed. As a result, the inner halo becomes less cuspy, while the outer halo 
retains its coherence. The disk forms within a gravitational field that is smoother 
and more extended than in the standard cosmological model.

The angular momentum distribution of the collapsing gas is similarly affected. 
Modes near the relaxation scale dominate the torque field, leading to a more 
coherent spin axis and a narrower distribution of specific angular momentum.
% ---------------------------------------------------------
\section{Derivation of the RDCC Rotation Curve}

The circular velocity of a test particle in the disk is given by
\begin{equation}
    v_{c}^{2}(r)
    = r\,\frac{\partial \Phi_{\mathrm{RDCC}}}{\partial r}.
\end{equation}

To evaluate this expression, we consider the RDCC-modified density profile of the 
halo. For a spherically symmetric halo, the potential satisfies
\begin{equation}
    \frac{\partial \Phi_{\mathrm{RDCC}}}{\partial r}
    = \frac{G M_{\mathrm{RDCC}}(r)}{r^{2}},
\end{equation}
where $M_{\mathrm{RDCC}}(r)$ is the enclosed mass.

The RDCC mass profile can be approximated as
\begin{equation}
    M_{\mathrm{RDCC}}(r)
    = M(r)
      \left[
        1 - \chi_{M}
        \left(\frac{1}{k_{\mathrm{rel}} r}\right)^{\alpha}
      \right],
\end{equation}
reflecting the reduced contribution of small-scale modes to the inner density.

Substituting this into the circular velocity yields
\begin{equation}
    v_{c,\mathrm{RDCC}}^{2}(r)
    = \frac{G M(r)}{r}
      \left[
        1 - \chi_{M}
        \left(\frac{1}{k_{\mathrm{rel}} r}\right)^{\alpha}
      \right].
\end{equation}

This expression captures the essential RDCC modification: the rotation curve rises 
more gently in the inner region and transitions smoothly to a flat profile at 
intermediate radii.
% ---------------------------------------------------------
\section{The Flatness of Rotation Curves in RDCC}

The flatness of rotation curves arises when the enclosed mass grows approximately 
linearly with radius. In RDCC, this behavior emerges naturally from the interplay 
between suppressed small-scale modes and coherent large-scale modes. The mass 
profile becomes
\begin{equation}
    M_{\mathrm{RDCC}}(r)
    \approx \rho_{0}\, r^{3}
    \left[
        1 - \chi_{\rho}
        \left(\frac{1}{k_{\mathrm{rel}} r}\right)^{\alpha}
    \right],
\end{equation}
which yields a circular velocity
\begin{equation}
    v_{c,\mathrm{RDCC}}(r)
    \approx \sqrt{4\pi G \rho_{0}}\, r
    \left[
        1 - \frac{\chi_{\rho}}{2}
        \left(\frac{1}{k_{\mathrm{rel}} r}\right)^{\alpha}
    \right].
\end{equation}

At radii $r \gg k_{\mathrm{rel}}^{-1}$, the correction term becomes negligible, and 
the rotation curve approaches a constant value. The flatness is therefore a direct 
consequence of the relaxation scale, which determines the radius at which the 
transition occurs.
% ---------------------------------------------------------
\section{Declining Rotation Curves and the Outer Halo}

At sufficiently large radii, the density profile of the halo declines more steeply, 
and the rotation curve begins to fall. In RDCC, this decline is smoother than in 
the standard cosmological model because the outer halo retains a higher degree of 
coherence. The circular velocity in the outer region can be approximated as
\begin{equation}
    v_{c,\mathrm{RDCC}}(r)
    \propto r^{-1/2}
    \left[
        1 - \chi_{\mathrm{out}}
        \left(\frac{1}{k_{\mathrm{rel}} r}\right)^{\alpha}
    \right].
\end{equation}

This expression predicts a gentle decline in the rotation curve, consistent with 
observations of extended HI disks.
% ---------------------------------------------------------
\section{Conclusion}

The dynamics of galactic disks in the Relaxation-Driven Cyclic Cosmology reflect 
the scale-dependent evolution of the gravitational potential. The relaxation 
mechanism modifies the density structure of halos, the distribution of angular 
momentum, and the gravitational field experienced by the disk. These effects 
produce rotation curves that rise gently in the inner region, remain flat over a 
broad range of radii, and decline smoothly in the outer halo. The present 
Companion establishes the theoretical foundation for future studies of disk 
stability, bar formation, and the interplay between baryons and dark matter in 
RDCC.

% ========================
% References
% ========================

\begin{thebibliography}{10}

\bibitem{Flagship} Lehmann, M. (2026). Relaxation-Driven Cyclic Cosmology (RDCC): A Minimal CPT-Symmetric Framework with a Single Parameter. Flagship Paper. Zenodo: \url{https://zenodo.org/records/19744958} (v43.0 or later).

\bibitem{Zenodo} The complete RDCC companion series (I--LVII) is available at the same Zenodo record above.

% Thematisch relevante Companions
\bibitem{CompXVI} Companion XVI: Boltzmann Code and Modified Hierarchy.
\bibitem{CompXXXI} Companion XXXI: RDCC Halo Model and Non-Linear Structure.
\bibitem{CompXXXII} Companion XXXII--XXXIX: Galaxy--Halo Connection, Cosmic Web, Lensing, RSD, ISW, tSZ, Rotation Curves, CGM.

% Wichtige externe Referenzen
\bibitem{Planck2020} Planck Collaboration (2020), Astron. Astrophys. 641, A6.
\bibitem{DESI2024} DESI Collaboration (2024/2025), arXiv: [aktuelle $f\sigma_8$/BAO-Paper].
\bibitem{JWST2024} Maiolino et al. (2024), Nature (JWST high-$z$ galaxies/quasars).
\bibitem{Kaiser1987} Kaiser, N. (1987), Mon. Not. R. Astron. Soc. 227, 1 (RSD).
\bibitem{Bartelmann2001} Bartelmann, M. \& Schneider, P. (2001), Phys. Rep. 340, 291 (Weak Lensing Review).
\bibitem{HuOkamoto2002} Hu, W. \& Okamoto, T. (2002), Astrophys. J. 574, 566 (CMB Lensing).
\bibitem{LewisChallinor2006} Lewis, A. \& Challinor, A. (2006), Phys. Rep. 429, 1 (CMB Lensing Review).
\bibitem{NFW1997} Navarro, J. F., Frenk, C. S. \& White, S. D. M. (1997), Astrophys. J. 490, 493 (Halo Profiles).

\end{thebibliography}

\end{document}
